\subsection{可积子丛}

本节中，X 表示一个 \(C^{r}\) 类流形，F 表示 T(X) 的 \(C^{s}\) 类向量子丛，其中 r、s 属于 \(N_{K}\)，且 \(s \leqslant r - 1\)。

9.3.1. 设 V 是 \(C^{k}\) 类流形（\(k \in N_{K}, k \leqslant r\)），且 f 是从 V 到 X 的 \(C^{k}\) 类态射。若 T(f) 将 T(V) 映入 F，则称 f 是 F 的积分。

X 的 \(C^{k}\) 类子流形 Z 称为 F 的积分子流形，如果嵌入 \(Z \rightarrow X\) 按上述意义是 F 的积分，即对所有 \(\mathrm{T}_{x}(\mathbf{Z}) \subset \mathrm{F}_{x}\) 都有 \(x \in Z\)。

设 \(x \in X\)。若 \(Z_1\) 和 \(Z_2\) 是包含 \(F\) 的 \(x\) 的两个积分子流形，且 \(T_x(Z_1) = F_x\)，则存在一个邻域 \(U\)，包含 \(x\)，使得 \(U \cap Z_2 \subset Z_1\)。若此外 \(T_x(Z_2) = F_x\)，则 \(Z_1\) 和 \(Z_2\) 在 \(x\) 处的芽相同。

9.3.2. 若存在 X 的 \(C^{s}\) 类叶状结构 Y，使得 \(F = T(X, Y)\)（9.2.8），则称 F 是可积的。此时该叶状结构是唯一的；称为 F 的积分叶状结构。\(V \to X\) 类态射 f: \(C^{s}\) 是 F 的积分的充要条件是 f 是从 V 到 Y 的态射。

例。--- 若 F 在每一点的秩为 1，则 F 可积，并在 X 上定义一个 1 维纯叶状结构。此外，假设 K = R，X 是分离的，且 F 有一个处处非零的向量场 \(\xi\) 作为标架。则对每个 \(x \in X\)，包含 \(x\) 的极大连通叶是以 \(\xi\) 为起点的 \(x\) 的最大积分弧的像（9.1.3）。

9.3.3（“可积性判据”）。以下条件等价：

\begin{enumerate}
\def\labelenumi{(\roman{enumi})}
\item
  T(X) 的向量子丛 F 是可积的。
\item
  对于每个 \(x \in X\)，存在一个 \(Z\)，它是 \(C^s\) 类 \(F\) 的积分子流形，使得 \(x \in Z\) 且 \(T_x(Z) = F_x\)。
\item
  对 X 的每个开集 U 以及属于 \(\xi\)（7.4.1）的向量场 \(\eta\) 和 \(\mathcal{S}_{\mathbf{F}}^{s}(\mathbf{U})\)，都有 \([\xi, \eta] \in \mathcal{S}_{\mathbf{F}}^{s-1}(\mathbf{U})\)。
\end{enumerate}

设 \((\xi_{i})_{i\in I}\) 是 F 的 \(C^{s}\) 类截面族，使得对每个 \(x\in X\)，\(\xi_{i}(x)\) 所成的集合是 Banach 空间 \(F_{x}\) 的总子集（EVT, I, § 2, no. 1）。则条件 (i)、(ii) 和 (iii) 等价于：

\begin{enumerate}
\def\labelenumi{(\roman{enumi})}
\setcounter{enumi}{3}
\tightlist
\item
  对 I 中每一对元素 \((i,j)\) 以及每个 \(x\in X\)，都有 \([\xi_i,\xi_j](x)\in F_x\)。
\end{enumerate}

当 I 有限且族 \((\xi_{i})\) 是 F 的一个标架时，前述条件仍等价于：

\begin{enumerate}
\def\labelenumi{(\alph{enumi})}
\setcounter{enumi}{21}
\tightlist
\item
  存在一个函数族 \((c_{ij}^{k})_{(i,j,k)\in I\times I\times I}\)，定义在 \(X\) 上、取值于 \(K\)，使得 \([\xi_i,\xi_j] = \sum_{k}c_{ij}^k\xi_k\)，对所有 \(i,j\)（属于 \(I\)）均成立。
\end{enumerate}

（若情形如此，则函数 \(c_{i,j}^{k}\) 属于 \(\mathbf{C}^{s-1}\) 类。）

9.3.4. 设 \(s' \in \mathbf{N}_{\mathbb{K}}\)，且 \(s' \leqslant s\)，令 \(\mathbf{F}'\) 为由 \(\mathbf{C}^{s}\) 通过弱化结构得到的 \(\mathbf{F}\) 类向量丛（8.7.1）。\(\mathbf{F}\) 可积的充要条件是 \(\mathbf{F}'\) 可积。在此情形，\(\mathbf{F}'\) 的积分叶状结构由 \(\mathbf{F}\) 的积分叶状结构通过弱化结构得到。

9.3.5. 设 E 是 Banach 空间，\(p\) 是一个 \(\geqslant 0\) 的整数。对每个 \(x \in X\)，记

\(N(p, E)_{x}\) 是 \(\text{Alt}^{p}(T_{x}(X); E)\) 的向量子空间，它由元素 \(u\) 组成，满足 \(u(v_{1}, \ldots, v_{p}) = 0\)，其中 \(v_{1}, \ldots, v_{p}\) 属于 \(F_{x}\)。空间 \(N(p, E)_{x}\)（其中 \(x \in X\)）构成一个 \(C^{s}\) 类向量子丛 \(\text{Alt}^{p}(T(X); E)\) 的纤维。记为 \(N(p, E)\)。若 \(F\) 可积，则有：

\begin{enumerate}
\def\labelenumi{(\roman{enumi})}
\setcounter{enumi}{5}
\tightlist
\item
  对 X 的每个开集 U 以及每个形式 \(\omega \in \mathcal{S}_{\mathrm{N}(p,\mathbb{E})}^s (\mathrm{U})\)，都有 \(d\omega \in \mathcal{S}_{\mathrm{N}(p + 1,\mathbb{E})}^{s - 1}(\mathrm{U})\)。
\end{enumerate}

反之，若 (vi) 对 \(p=1\) 以及每个 Banach 空间 E 都成立，则丛 F 可积；当 K = R 或 C 时，甚至只需验证 \(p=1\) 且 E = K 时的 (vi)。

假设 T(X)/F 的对偶具有一个标架 \((\omega_{1},\ldots,\omega_{n})\)。则 F 的可积性等价于下列条件：

\begin{enumerate}
\def\labelenumi{(\roman{enumi})}
\setcounter{enumi}{6}
\tightlist
\item
  \[
  d \omega_ {i} \wedge \omega_ {1} \wedge \dots \wedge \omega_ {n} = 0 \quad \text {for} \quad 1 \leqslant i \leqslant n.
  \]
\end{enumerate}

若此外存在 T(X) 的 \(C^{s}\) 类向量子丛 G，使得 T(X) 是 F 与 G 的直和，则条件 (vii) 等价于：

\begin{enumerate}
\def\labelenumi{(\roman{enumi})}
\setcounter{enumi}{7}
\tightlist
\item
  存在 X 上次数为 1、属于 \(\alpha_{i}^{j}\) 类的微分形式 \(i, j\)（\(\mathbf{C}^{s-1}\) 属于 I），使得对所有 \(d\omega_{i} = \sum_{j} \alpha_{i}^{j} \wedge \omega_{j}\)，有 \(i \in \mathbf{I}\)。
\end{enumerate}

9.3.6. 设 L 是 Banach 空间，且 \(\omega\) 是 X 上取值于 L、属于 C \(^{s}\) 类的次数为 1 的微分形式。作下列两个假设：

\begin{enumerate}
\def\labelenumi{\alph{enumi})}
\item
  对所有 \(x \in \mathbf{X}\)，\(\omega_x\) 是从 \(\mathrm{T}_x(\mathbf{X})\) 到 \(\mathbf{L}\) 的满同态，且 \(\omega_x\) 的核是 \(\mathrm{F}_x\)；
\item
  存在 T(X) 的一个向量子丛 G，使得 T(X) 是 F 与 G 的直和。
\end{enumerate}

则 F 的可积性等价于：

\begin{enumerate}
\def\labelenumi{(\roman{enumi})}
\setcounter{enumi}{8}
\tightlist
\item
  存在 X 上取值于 \(\alpha\) 的次数为 1 的微分形式 \(\operatorname{End}(\mathbf{L})\)，使得 \(d\omega = \alpha \wedge \omega\)，其中外积由典范配对 \(\operatorname{End}(\mathbf{L}) \times \mathbf{L} \to \mathbf{L}\) 定义（7.8.2 和 8.3.2）。
\end{enumerate}

9.3.7（“全微分方程”）。假设 X 是流形 A 和 B 的乘积，且属于 \(C^{r}\) 类；记 \(p_{1}: X \to A\) 和 \(p_{2}: X \to B\) 为两个投影。设 f 是从 \(C^{s}\) 到 \(p_{1}^{*}T(A)\) 的 \(p_{2}^{*}T(B)\) 类态射。映射 \(f_{(a,b)}: T_{a}(A) \to T_{b}(B)\) 的图是 \(C^{s}\) 的一个 \(T(X)\) 类向量子丛的纤维；记为 \(F^{f}\)。

设 A' 是 A 的开子集，且 \(\varphi: \mathbf{A}' \to \mathbf{B}\) 是 \(\mathbf{C}^k\) 类态射（\(k \in \mathbb{N}_{\mathbb{K}}\)，\(k \leqslant r\)）。称 \(\varphi\) 为 \(f\) 的积分，如果对所有 \(a \in \mathbf{A}'\) 都有 \(\mathrm{T}_a(\varphi) = f_{a,\varphi(a)}\)；这等价于说，映射 \(a \mapsto (a, \varphi(a))\)（从 \(\mathbf{A}'\) 到 X）是 \(\mathbf{F}^f\) 的积分（9.3.1）。这样的映射属于 \(\mathbf{C}^{s+1}\) 类。若 \(\varphi_1\) 和 \(\varphi_2\) 是 \(f\) 的两个积分，并在 \(a \in \mathbf{A}\) 的一点处取相同值，则它们在 \(a\) 的某个邻域内一致。

更一般地，设 Z 是 \(C^{k}\) 类流形，A' 是 A 的开子集，\(a \in A'\)，且 \(\Phi_{1}, \Phi_{2}\) 是从 \(C^{k}\) 到 B 的 \(Z \times A'\) 类态射。

若 \(\Phi_{1}\) 和 \(\Phi_{2}\) 在 \(Z \times \{a\}\) 上一致，并且对所有 \(z \in Z\)，态射

\[
a \mapsto \Phi_ {1} (z, a) \quad \text {and} \quad a \mapsto \Phi_ {2} (z, a)
\]

是 \(f\) 的积分，则 \(\Phi_{1}\) 和 \(\Phi_{2}\) 在 \(Z \times \{a\}\) 的某个邻域内一致。

假设 \(F^{f}\) 可积。设 Z 是 \(C^{k}\) 类流形（\(k \in \mathbf{N}_{\mathbb{K}}, k \leqslant s\)），\((z_{0}, a_{0})\) 是 \(Z \times A\) 中的一点，且 \(\rho\) 是从 Z 到 B 的 \(C^{k}\) 类态射。存在一个开邻域 \(Z' \times A'\)，它是 \((z_{0}, a_{0})\) 在 \(Z \times A\) 中的邻域，以及一个 \(\Phi: Z' \times A' \to B\) 类态射 \(C^{k}\)，使得对所有 \(z \in Z'\)，映射 \(a \mapsto \Phi(z, a)\) 从 \(A'\) 到 B，以 f 为其切映射，并在点 \(\rho(z)\) 处取值 \(a_{0}\)。

9.3.8. 保持 9.3.7 的假设和记号，并假设 A（分别为 B）是 Banach 空间 E（分别为 M）的开子流形。此时，映射 f 可视为从 X = A × B 到 Banach 空间 \(^{s}\) 的 C \(\mathcal{L}(E; M)\) 类态射。将 f 的第一（分别第二）个偏导数记为 D \(_{1}\) f（分别 D \(_{2}\) f）（1.6.2）；它是从 X 到 \(^{s-1}\)（分别到 \(\mathcal{L}(E; \mathcal{L}(E; M))\)）的 C \(\mathcal{L}(M; \mathcal{L}(E; M))\) 类态射，该空间以显然的方式与 \(\mathcal{L}_{2}(E; M)\)（分别与 \(\mathcal{L}(M, E; M)\)）等同。F 可积的充要条件是，对每个 x ∈ X，双线性映射

\[
\Delta_ {x}: (h _ {1}, h _ {2}) \mapsto \mathrm{D} _ {1} f (x) (h _ {1}, h _ {2}) + \mathrm{D} _ {2} f (x) (f (x) h _ {1}, h _ {2})
\]

从 \(E \times E\) 到 M 的双线性映射是对称的。在此条件下，若 \(\varphi\) 是定义在 A 的开子集 A' 上的 f 的积分，则 \(\varphi\) 在 A' 的一点 \(a\) 处的二阶导数为 \(\Delta(a,\varphi(a))\)。

若 \(E = K^n\)，并记 \((x^1, \ldots, x^n)\) 为 \(K^n\) 上的坐标函数，则映射 \(f\) 由从 \((f_1, \ldots, f_n)\) 到 \(A \times B\) 的映射族 \(M\) 定义，且可积性条件写为：

\[
\frac {\partial f _ {i}}{\partial x ^ {j}} + (\mathbf {D} _ {2} f _ {i}). f _ {j} = \frac {\partial f _ {j}}{\partial x ^ {i}} + (\mathbf {D} _ {2} f _ {j}). f _ {i}
\]

对所有整数 \(i, j\)（它们属于 \([1, n]\)），A 的开子集 A' 到 B 的映射 \(\varphi\) 是 \(f\) 的积分，当且仅当有

\[
\frac {\partial \varphi}{\partial x ^ {i}} = f _ {i} (x, \varphi (x)) \quad \text {for all} \quad x \in \mathbf {A} ^ {\prime} \text { and all } i \in [ 1, n ],
\]

换言之，有

\[
d \varphi = \sum_ {1 \leqslant i \leqslant n} f _ {i} (x, \varphi (x)) d x ^ {i}.
\]

